% =====================================================================
% FONTES DEPENDENTES EM MALHAS --- NIVEL 3
% Questoes integradas ao bloco de Aprofundamento do Capitulo 9.
% Correntes de malha desenhadas no sentido HORARIO.
% =====================================================================

\begin{questao}[3]{}
A fonte dependente vale $4i_x$ volts, com $i_x=I_1-I_2$ no resistor central.
Determine as correntes de malha.

\circuitoq{%
\begin{circuitikz}[american,scale=.92,transform shape,line width=.8pt]
\draw (0,0) to[\fonteV,l^={$\SI{24}{\volt}$}] (0,3.2)
      to[R,l={$\SI{4}{\ohm}$}] (3.2,3.2) coordinate(t);
\draw (t) to[R,l={$\SI{4}{\ohm}$},i>^={$i_x$}] (3.2,0) coordinate(b) -- (0,0);
\draw (t) to[R,l={$\SI{8}{\ohm}$}] (6.4,3.2)
      to[cvsource,l={$4i_x$}] (6.4,0) -- (b);
\draw[->,loopcol,line width=1.2pt] (1.75,1.00) arc[start angle=0,end angle=-310,radius=.46cm];
\node[loopcol,fill=white,inner sep=1pt] at (1.75,1.00){\footnotesize$I_1$};
\draw[->,loopcol,line width=1.2pt] (4.75,1.00) arc[start angle=0,end angle=-310,radius=.46cm];
\node[loopcol,fill=white,inner sep=1pt] at (4.75,1.00){\footnotesize$I_2$};
\end{circuitikz}}

\resposta{$I_1=\SI{4}{\ampere}$; $I_2=\SI{2}{\ampere}$}
\resolucao{A malha 1 fornece $8I_1-4I_2=24$. Na malha 2,
\[
-4I_1+12I_2=4(I_1-I_2),
\]
ou $-8I_1+16I_2=0$. Assim $I_1=2I_2$ e, pela primeira equação,
$I_1=\SI{4}{\ampere}$ e $I_2=\SI{2}{\ampere}$.}
\end{questao}

\begin{questao}[3]{}
A tensão de controle $v_x$ é medida no resistor de $\SI{4}{\ohm}$ exclusivo da
malha 1. A fonte dependente vale $0{,}5v_x$. Determine $I_1$ e $I_2$.

\circuitoq{%
\begin{circuitikz}[american,scale=.92,transform shape,line width=.8pt]
\draw (0,0) to[\fonteV,l^={$\SI{24}{\volt}$}] (0,3.2)
      to[R,l_={$\SI{4}{\ohm}$},v^>={$v_x$}] (3.2,3.2) coordinate(t);
\draw (t) to[R,l={$\SI{4}{\ohm}$}] (3.2,0) coordinate(b) -- (0,0);
\draw (t) to[R,l={$\SI{4}{\ohm}$}] (6.4,3.2)
      to[cvsource,l={$0{,}5v_x$}] (6.4,0) -- (b);
\draw[->,loopcol,line width=1.2pt] (1.75,1.00) arc[start angle=0,end angle=-310,radius=.46cm];
\node[loopcol,fill=white,inner sep=1pt] at (1.75,1.00){\footnotesize$I_1$};
\draw[->,loopcol,line width=1.2pt] (4.75,1.00) arc[start angle=0,end angle=-310,radius=.46cm];
\node[loopcol,fill=white,inner sep=1pt] at (4.75,1.00){\footnotesize$I_2$};
\end{circuitikz}}

\resposta{$I_1=\SI{4.8}{\ampere}$; $I_2=\SI{3.6}{\ampere}$}
\resolucao{Como $v_x=4I_1$, a fonte dependente vale $2I_1$ volts. Logo,
\[
8I_1-4I_2=24,
\qquad
-4I_1+8I_2=2I_1.
\]
Da segunda, $I_2=0{,}75I_1$. Substituindo na primeira, obtêm-se
$I_1=\SI{4.8}{\ampere}$ e $I_2=\SI{3.6}{\ampere}$.}
\end{questao}

\begin{questao}[3]{}
A fonte de corrente dependente está no ramo comum e fornece corrente de baixo
para cima igual a $0{,}1v_x$. A tensão $v_x$ está indicada no resistor da malha
1. Determine $I_1$ e $I_2$ por supermalha.

\circuitoq{%
\begin{circuitikz}[american,scale=.92,transform shape,line width=.8pt]
\draw (0,0) to[\fonteV,l^={$\SI{40}{\volt}$}] (0,3.2)
      to[R,l_={$\SI{5}{\ohm}$},v^>={$v_x$}] (3.2,3.2) coordinate(t);
\draw (t) to[cisource,l={$0{,}1v_x$},invert] (3.2,0) coordinate(b) -- (0,0);
\draw (t) to[R,l={$\SI{5}{\ohm}$}] (6.4,3.2) -- (6.4,0) -- (b);
\draw[->,loopcol,line width=1.2pt] (1.75,1.00) arc[start angle=0,end angle=-310,radius=.46cm];
\node[loopcol,fill=white,inner sep=1pt] at (1.75,1.00){\footnotesize$I_1$};
\draw[->,loopcol,line width=1.2pt] (4.75,1.00) arc[start angle=0,end angle=-310,radius=.46cm];
\node[loopcol,fill=white,inner sep=1pt] at (4.75,1.00){\footnotesize$I_2$};
\draw[dashed,laranja,line width=.9pt,rounded corners=7pt]
      (-.55,-.5) rectangle (6.95,4.10);
\node[laranja,font=\footnotesize\sffamily,fill=white,inner sep=1.5pt]
      at (3.2,-.80) {supermalha};
\end{circuitikz}}

\resposta{$I_1=\SI{3.2}{\ampere}$; $I_2=\SI{4.8}{\ampere}$}
\resolucao{Como $v_x=5I_1$, a restrição imposta pela fonte é
\[
I_2-I_1=0{,}1(5I_1)=0{,}5I_1,
\]
portanto $I_2=1{,}5I_1$. A LKT na supermalha externa é
$40=5I_1+5I_2$. Substituindo, $I_1=\SI{3.2}{\ampere}$ e
$I_2=\SI{4.8}{\ampere}$.}
\end{questao}

\begin{questao}[3]{}
A fonte de corrente dependente do ramo comum fornece, de baixo para cima,
$2i_x$, sendo $i_x=I_1$ no resistor esquerdo. Determine as correntes de malha
por supermalha.

\circuitoq{%
\begin{circuitikz}[american,scale=.92,transform shape,line width=.8pt]
\draw (0,0) to[\fonteV,l^={$\SI{40}{\volt}$}] (0,3.2)
      to[R,l_={$\SI{5}{\ohm}$},i>^={$i_x$}] (3.2,3.2) coordinate(t);
\draw (t) to[cisource,l={$2i_x$},invert] (3.2,0) coordinate(b) -- (0,0);
\draw (t) to[R,l={$\SI{5}{\ohm}$}] (6.4,3.2) -- (6.4,0) -- (b);
\draw[->,loopcol,line width=1.2pt] (1.75,1.00) arc[start angle=0,end angle=-310,radius=.46cm];
\node[loopcol,fill=white,inner sep=1pt] at (1.75,1.00){\footnotesize$I_1$};
\draw[->,loopcol,line width=1.2pt] (4.75,1.00) arc[start angle=0,end angle=-310,radius=.46cm];
\node[loopcol,fill=white,inner sep=1pt] at (4.75,1.00){\footnotesize$I_2$};
\draw[dashed,laranja,line width=.9pt,rounded corners=7pt]
      (-.55,-.5) rectangle (6.95,4.10);
\node[laranja,font=\footnotesize\sffamily,fill=white,inner sep=1.5pt]
      at (3.2,-.80) {supermalha};
\end{circuitikz}}

\resposta{$I_1=\SI{2}{\ampere}$; $I_2=\SI{6}{\ampere}$}
\resolucao{A restrição é $I_2-I_1=2I_1$, portanto $I_2=3I_1$. A LKT na
supermalha é $40=5I_1+5I_2$. Assim $40=20I_1$, resultando
$I_1=\SI{2}{\ampere}$ e $I_2=\SI{6}{\ampere}$.}
\end{questao}
